\documentclass[11pt,a4paper]{article}

% --- Codificación y Lenguaje para pdfLaTeX ---
\usepackage[utf8]{inputenc}
\usepackage[T1]{fontenc}
\usepackage[spanish,es-tabla]{babel}

% --- Tipografía estándar ---
\usepackage{lmodern}

% --- Geometría de Página ---
\usepackage[margin=2.5cm]{geometry}

% --- Matemáticas ---
\usepackage{amsmath,amssymb,amsfonts}

% --- Tablas y Gráficos ---
\usepackage{graphicx}
\usepackage{booktabs}
\usepackage{tabularx}
\usepackage{subcaption}
\usepackage{xcolor}

% --- Algoritmos y Pseudocódigo ---
\usepackage{algorithm}
\usepackage{algpseudocode}
% Traducción de palabras clave de algoritmos a español
\makeatletter
\renewcommand{\algorithmicrequire}{\textbf{Entrada:}}
\renewcommand{\algorithmicensure}{\textbf{Salida:}}
\makeatother

% --- Código Fuente ---
\usepackage{listings}
\lstset{
    basicstyle=\ttfamily\small,
    breaklines=true,
    frame=single,
    numbers=left,
    numberstyle=\tiny,
    showstringspaces=false
}

% --- Hipervínculos ---
\usepackage{hyperref}
\hypersetup{
    colorlinks=true,
    linkcolor=black,
    filecolor=black,
    urlcolor=blue,
    citecolor=blue
}

% --- Encabezados y Pies de Página ---
\usepackage{fancyhdr}
\setlength{\headheight}{14pt}
\pagestyle{fancy}
\fancyhf{}
\fancyhead[L]{TB024 - Teoría de Algoritmos (FIUBA)}
\fancyhead[R]{TP I - 2C 2026 - Grupo 17}
\fancyfoot[C]{\thepage}

\begin{document}

% ==========================================
% PORTADA / CARÁTULA
% ==========================================
\begin{titlepage}
    \centering
    \vspace*{1cm}

    {\Huge \textbf{Universidad de Buenos Aires}}\\[0.3cm]
    {\LARGE Facultad de Ingeniería}\\[1.5cm]

    {\Large \textbf{TB024 -- Teoría de Algoritmos}}\\[0.2cm]
    {\large Cátedra: Echevarria}\\[0.2cm]
    {\large 2do Cuatrimestre de 2026}\\[2cm]

    \rule{\linewidth}{0.5mm}\\[0.4cm]
    {\Huge \textbf{Trabajo Práctico I}}\\[0.3cm]
    {\Large Problema de la mochila}\\[0.4cm]
    \rule{\linewidth}{0.5mm}\\[2.5cm]

    \textbf{\large Grupo N$^\circ$ 17}\\[0.8cm]

    \begin{tabular}{llr}
        \toprule
        \textbf{Apellido y Nombre} & \textbf{Padrón} \\
        \midrule
        Barberis, Juan Celestino & 105147 \\
        Conte, Valentino & 112402 \\
        Gorge, Julián Alfredo & 104286 \\
        Laruse De Moraes Fernandes, Giovanna & 113100 \\
        \bottomrule
    \end{tabular}

    \vfill
\end{titlepage}

% ==========================================
% ÍNDICE GENERAL
% ==========================================
\tableofcontents
\newpage

% ==========================================
% SECCIÓN 0: ENTORNO Y REPRODUCIBILIDAD
% ==========================================
\section{Instrucciones de ejecución y entorno}
\label{sec:entorno}

\subsection{Lenguaje y dependencias}
\begin{itemize}
    \item \textbf{Lenguaje de programación:} Python 3.
    \item \textbf{Bibliotecas requeridas:} \texttt{pulp}, \texttt{matplotlib}.
\end{itemize}

\subsection{Instrucciones de compilación y ejecución}
Las instrucciones para correr el proyecto y las pruebas están en el README.md del proyecto.

\newpage

\section{Algoritmo 1: Fuerza bruta}

\subsection{Supuestos}
\begin{itemize}
    \item \textbf{Limitación de tamaño:} El algoritmo evalúa exhaustivamente las $2^N$ combinaciones posibles sin aplicar ningún tipo de filtro prematuro, lo que lo vuelve inviable en tiempo para mochilas de tamaño mediano o grande ($N > 25$).
    \item \textbf{Independencia del orden:} El orden inicial de los objetos no afecta el resultado final, ni tampoco el tiempo de ejecución, ya que el algoritmo siempre explorará la totalidad del espacio de soluciones.
\end{itemize}

\subsection{Diseño}
Para que el algoritmo sea estrictamente de fuerza bruta pura, debe generar todas las combinaciones posibles y recién al final evaluar su factibilidad (peso) y optimalidad (beneficio).

\textbf{Estructuras de datos utilizadas:}
\begin{itemize}
    \item \textbf{Arreglo/Lista de tuplas:} Para almacenar los objetos \texttt{(peso, beneficio)}.
\end{itemize}

\textbf{Pseudocódigo:}
\begin{lstlisting}[language=Python]
Funcion FuerzaBruta(capacidad, objetos):
    mejor_beneficio = 0
    mejor_combinacion = []

    Para cada tamaño r desde 0 hasta longitud(objetos):
        Para cada combinacion en GenerarCombinaciones(objetos, r):
            peso_total = SumarPesos(combinacion)

            Si peso_total <= capacidad:
                ben_total = SumarBeneficios(combinacion)
                Si ben_total > mejor_beneficio:
                    mejor_beneficio = ben_total
                    mejor_combinacion = combinacion

    Retornar (mejor_combinacion, mejor_beneficio)
\end{lstlisting}

\subsection{Seguimiento}
Ejemplo reducido: Capacidad $W = 5$. Objetos: \texttt{O1(peso:3, ben:4)}, \texttt{O2(peso:2, ben:3)}.

\begin{table}[h!]
    \centering
    \begin{tabular}{@{}lccc@{}}
        \toprule
        \textbf{Combinación generada} & \textbf{Peso total} & \textbf{¿Es factible? ($\le 5$)} & \textbf{Beneficio total} \\ \midrule
        \texttt{[]} (Vacía) & 0 & Sí & 0 \\
        \texttt{[O1]} & 3 & Sí & 4 \\
        \texttt{[O2]} & 2 & Sí & 3 \\
        \texttt{[O1, O2]} & 5 & Sí & \textbf{7 (Óptimo)} \\ \bottomrule
    \end{tabular}
\end{table}
Como se observa, Fuerza Bruta pura evalúa las $2^2 = 4$ combinaciones posibles antes de descartar aquellas que exceden el peso, garantizando hallar el máximo global.

\subsection{Complejidad}
\begin{itemize}
    \item \textbf{Temporal:} $\mathcal{O}(2^N)$. Genera obligatoriamente el conjunto potencia completo de los $N$ objetos.
    \item \textbf{Espacial:} $\mathcal{O}(N)$ en recursión, o $\mathcal{O}(1)$ adicional si las combinaciones se generan iterativamente, excluyendo el almacenamiento de los datos.
\end{itemize}

\newpage
\section{Algoritmo 2: Backtracking}

\subsection{Supuestos}
\begin{itemize}
    \item \textbf{Ordenamiento previo:} Asume que la lista de objetos se ordena de mayor a menor según la relación \textit{beneficio/peso} antes de iniciar la recursión. Esto es crucial para alcanzar rápidamente buenas soluciones que fortalezcan la poda por cota superior.
\end{itemize}

\subsection{Podas utilizadas}
Para evitar el comportamiento exponencial de Fuerza Bruta, este algoritmo implementa dos cortes fundamentales:
\begin{itemize}
    \item \textbf{Poda por Factibilidad:} Antes de profundizar en una rama incluyendo un objeto, se verifica que \texttt{peso\_acumulado + peso\_actual <= capacidad}. Si se excede, esa rama de exploración se corta inmediatamente.
    \item \textbf{Poda por Optimalidad (Cota Superior):} Mediante una relajación lineal (Mochila Fraccionaria), se proyecta el beneficio máximo posible a futuro (\texttt{cota}). Si esta cota es menor o igual al \texttt{mejor\_beneficio} ya descubierto, la rama se poda por completo.
\end{itemize}

\subsection{Diseño}
\textbf{Pseudocódigo:}
\begin{lstlisting}[language=Python]
    Funcion Explorar(indice, peso_acum, ben_acum):
        Si ben_acum > mejor_beneficio: mejor_beneficio = ben_acum
        Si indice == cantidad_objetos: Retornar

        // PODA 1
        cota = CalcularCotaSuperior(indice, peso_acum, ben_acum, ...)
        Si cota <= mejor_beneficio:
            Retornar

        // PODA 2
        Si peso_acum + peso_actual <= capacidad:
            Explorar(indice + 1, peso_acum + peso_actual, ben_acum + ben_actual)

        // Rama que excluye el objeto actual
        Explorar(indice + 1, peso_acum, ben_acum)

    Explorar(0, 0, 0)
\end{lstlisting}

\subsection{Seguimiento}
Para demostrar la efectividad de la poda por optimalidad,analizo un caso donde se utilice
Sea $W = 10$. Objetos (ya ordenados): \texttt{O1(p:4, b:40)}, \texttt{O2(p:7, b:42)}, \texttt{O3(p:5, b:25)}, \texttt{O4(p:3, b:12)}.
\begin{center}
    \begin{tabular}{@{}p{3.5cm}p{2cm}p{5.5cm}p{3cm}@{}}
        \toprule
        \textbf{Acción en la Rama} & \textbf{Estado} & \textbf{Evaluación / Cálculo de Cota} & \textbf{Resultado} \\ \midrule
        1. Incluye O1 & P=4, B=40 & Intenta incluir O2: $P=11>10$. & \textbf{Poda Factibilidad}. \\
        2. Omite O2, incluye O3 & P=9, B=65 & Mejor beneficio actualizado a 65. & \texttt{mejor\_ben} = 65 \\
        3. \textbf{Backtrack: Excluye O1} & P=0, B=0 & Cota tomando O2 y fracción de O3: $42 + (25 \times \frac{3}{5}) = 57$. & $57 \le 65$. \textbf{Poda Optimalidad}. \\ \bottomrule
        \bottomrule
    \end{tabular}
\end{center}





\subsection{Complejidad}
\begin{itemize}
    \item \textbf{Temporal (Peor Caso):} $\mathcal{O}(2^N)$. En instancias donde las podas no logran cortar ramas, el árbol binario se explora en su totalidad.
    \item \textbf{Temporal (Caso Promedio):} Menor a $\mathcal{O}(2^N)$. Las podas evitan explorar la mayoría del espacio de soluciones. Requiere un tiempo base de $\mathcal{O}(N \log N)$ para el ordenamiento previo.
\end{itemize}

\subsection{Análisis de Ejecución y Conclusión General}

{Escalabilidad y Crecimiento Empírico}
Se realizaron ejecuciones para medir el tiempo real y el comportamiento asintótico a medida que crecía el tamaño del problema ($N$). Para el análisis, se utilizaron progresiones de $N=5, 10, 15, 20, 25$, y conjuntos grandes extraídos de los archivos \texttt{mochila50} y \texttt{mochila90}.

El análisis de ejecución empírico arroja resultados contundentes sobre el diseño de los algoritmos:

\begin{itemize}
    \item \textbf{Fuerza Bruta:} Como se esperaba teóricamente, el algoritmo de Fuerza Bruta pura mostró tiempos de ejecución que crecieron de forma exponencial ($\approx 2^N$). A partir de $N=25$, el algoritmo se vuelve computacionalmente inviable en tiempos prácticos, debido a que carece de mecanismos para descartar combinaciones que fallan y se ve forzado a explorar millones de estados inútiles.

    \item \textbf{Backtracking:} Si bien el Backtracking comparte la misma complejidad teórica de peor caso asintótico ($\mathcal{O}(2^N)$), los resultados de las ejecuciones reales demostraron que el algoritmo logra procesar exitosamente sets de datos muy superiores (como $N=50$ o $N=90$) en fracciones de segundo.
\end{itemize}

\subsection{Conclusión}
La diferencia radical en los tiempos de ejecución observados no radica en la complejidad teórica, sino en la arquitectura del recorrido. Mientras que la Fuerza Bruta evalúa a ciegas todo el espacio de soluciones, el Backtracking utiliza la Cota Fraccionaria y restricciones de dominio para podar ramas en etapas muy tempranas.

\begin{center}
    \includegraphics[width=0.9\textwidth]{crecimiento_fb_vs_bt.png}
\end{center}
% ==========================================
% PROBLEMA 2: GREEDY
% ==========================================
\section{Problema 2: Algoritmo Greedy (Garantía de Calidad 1/2)}
\label{sec:problema2}

\subsection{Supuestos y Limitaciones}
\textbf{Supuestos:} Se considera el problema de la mochila 0/1, donde cada
objeto $i$ posee un peso $w_i$ y un beneficio $v_i$, y puede ser incluido
completamente en la mochila o excluido, sin posibilidad de fraccionarlo.
La mochila posee una capacidad máxima $W$.

Para aplicar el algoritmo con garantía $1/2$, se consideran únicamente los
objetos cuyo peso es menor o igual a $W$, es decir, aquellos que pueden entrar
individualmente en la mochila. Luego, estos objetos se ordenan de forma
decreciente según la relación beneficio/peso $v_i/w_i$.

\textbf{Limitaciones:} El algoritmo es aproximado y, por lo tanto, no garantiza
obtener la solución óptima. La garantía $1/2$ establece que el beneficio de la
solución obtenida será al menos la mitad del beneficio de una solución óptima.

\subsection{Diseño y estructuras de datos}

\textbf{Estructuras de datos utilizadas:}

\begin{itemize}
    \item \textbf{Lista de tuplas:} para almacenar los objetos como pares
    (peso, beneficio).

    \item \textbf{Lista de objetos seleccionados:} almacena los objetos
    incorporados a la solución greedy.
\end{itemize}

El algoritmo comienza descartando los objetos cuyo peso supera la capacidad
$W$. Luego, ordena los objetos restantes de forma decreciente según su relación
beneficio/peso $v_i/w_i$.

Se recorren los objetos en ese orden y se incorporan mientras exista capacidad.
El primer objeto que no puede ser agregado se denomina \textit{elemento crítico}.
Se compara entonces el beneficio acumulado de los objetos seleccionados
($P_1$) con el beneficio del elemento crítico ($P_2$), y se devuelve la
alternativa de mayor beneficio.

\paragraph{Pseudocódigo:}
\begin{lstlisting}
Funcion Mochila(W, objetos):
    objetosValidos = objetos cuyo peso <= W
    ordenar objetosValidos de mayor a menor segun vi/wi

    seleccionados = lista vacia
    pesoAcumulado = 0
    beneficioAcumulado = 0
    elementoCritico = ninguno

    para cada objeto en objetosValidos:
        si pesoAcumulado + objeto.peso <= W:
            agregar objeto a seleccionados
            pesoAcumulado = pesoAcumulado + objeto.peso
            beneficioAcumulado = beneficioAcumulado + objeto.beneficio
        sino:
            elementoCritico = objeto
            cortar

    si elementoCritico = ninguno:
        devolver seleccionados, beneficioAcumulado

    si beneficioAcumulado >= elementoCritico.beneficio:
        devolver seleccionados, beneficioAcumulado
    sino:
        devolver [elementoCritico], elementoCritico.beneficio
\end{lstlisting}

\subsection{Seguimiento con set reducido}

Se considera una mochila de capacidad $W=10$ y los siguientes objetos,
ya ordenados de forma decreciente según su relación beneficio/peso:

\begin{table}[h]
    \centering
    \begin{tabular}{cccc}
        \toprule
        \textbf{Objeto} & \textbf{Peso ($w_i$)} & \textbf{Beneficio ($v_i$)} & \textbf{$v_i/w_i$} \\
        \midrule
        A & 4 & 20 & 5 \\
        B & 5 & 20 & 4 \\
        C & 3 & 9 & 3 \\
        \bottomrule
    \end{tabular}
    \caption{Set reducido para el seguimiento del algoritmo Greedy.}
\end{table}

Inicialmente, la lista de seleccionados está vacía, el peso acumulado es 0
y el beneficio acumulado es 0.

\begin{enumerate}
    \item Se analiza el objeto A. Como $0+4 \leq 10$, se incorpora.
    El peso acumulado pasa a ser 4 y el beneficio acumulado 20.

    \item Se analiza el objeto B. Como $4+5 \leq 10$, se incorpora.
    El peso acumulado pasa a ser 9 y el beneficio acumulado 40.

    \item Se analiza el objeto C. Como $9+3 > 10$, no puede incorporarse
    y se convierte en el elemento crítico.
\end{enumerate}

Por lo tanto, la solución construida tiene beneficio $P_1=40$, mientras que
el elemento crítico tiene beneficio $P_2=9$. Como $P_1 > P_2$, el algoritmo
devuelve los objetos A y B, con peso total 9 y beneficio total 40.

\subsection{Demostración y caso de test para garantía 1/2}

Sean $P_1$ el beneficio de los objetos incorporados por el algoritmo greedy
antes de encontrar el elemento crítico, y $P_2$ el beneficio de dicho elemento.

Como los objetos se encuentran ordenados de forma decreciente según
$v_i/w_i$, la solución fraccionaria que toma los objetos anteriores al crítico
y, si fuera necesario, una fracción de este último constituye una cota superior
para la solución óptima de la mochila 0/1.

Por lo tanto:

\[
OPT \leq P_1 + P_2
\]

Por otra parte, el algoritmo devuelve la alternativa de mayor beneficio entre
$P_1$ y $P_2$. Para dos valores no negativos se cumple:

\[
\max(P_1,P_2) \geq \frac{P_1+P_2}{2}
\]

Combinando ambas desigualdades:

\[
\max(P_1,P_2)
\geq \frac{P_1+P_2}{2}
\geq \frac{OPT}{2}
\]

Por lo tanto, la solución devuelta por el algoritmo posee un beneficio de al
menos la mitad del beneficio de la solución óptima, obteniéndose así una
garantía de calidad $1/2$.

\subsubsection{Set de datos diseñado ad-hoc}

Para mostrar la importancia de comparar la solución greedy con el elemento
crítico, se considera una mochila de capacidad $W=10$:

\begin{table}[h]
    \centering
    \begin{tabular}{cccc}
        \toprule
        \textbf{Objeto} & \textbf{Peso ($w_i$)} & \textbf{Beneficio ($v_i$)} & \textbf{$v_i/w_i$} \\
        \midrule
        A & 2 & 20 & 10 \\
        B & 9 & 81 & 9 \\
        \bottomrule
    \end{tabular}
    \caption{Caso de prueba donde resulta conveniente elegir el elemento crítico.}
\end{table}

El algoritmo incorpora primero el objeto A, obteniendo peso 2 y beneficio
$P_1=20$. Luego intenta incorporar B, pero $2+9>10$, por lo que B se convierte
en el elemento crítico.

El beneficio del elemento crítico es $P_2=81$. Como $P_2>P_1$, el algoritmo
descarta la solución formada por A y devuelve únicamente B, obteniendo un
beneficio de 81.

\subsection{Análisis teórico de complejidad}

El algoritmo realiza las siguientes operaciones principales:

\begin{itemize}
    \item Filtrar los objetos cuyo peso supera la capacidad $W$: $O(n)$.
    \item Ordenar los objetos válidos según la relación $v_i/w_i$: $O(n \log n)$.
    \item Recorrer los objetos ordenados hasta encontrar el elemento crítico:
    $O(n)$ en el peor caso.
    \item Comparar los beneficios $P_1$ y $P_2$: $O(1)$.
\end{itemize}

Por lo tanto, la complejidad temporal total es:

\[
O(n) + O(n \log n) + O(n) + O(1) = O(n \log n)
\]

El ordenamiento es la operación dominante, por lo que la complejidad temporal
final del algoritmo es $O(n \log n)$.

En cuanto a la complejidad espacial, se almacenan la lista de objetos válidos,
la lista ordenada y los objetos seleccionados. En el peor caso estas estructuras
pueden contener una cantidad lineal de elementos, por lo que la complejidad
espacial es $O(n)$.

\subsection{Resultados experimentales y ajuste curva teórica}

Para analizar empíricamente el comportamiento temporal del algoritmo Greedy,
se utilizaron sets de datos de distintos tamaños generados con el mismo criterio:
capacidad $W=50n$, pesos entre 1 y 200 y beneficios entre 1 y 1000.

Para cada tamaño se ejecutó el algoritmo 100 veces y se tomó el tiempo promedio,
con el objetivo de reducir el efecto de pequeñas variaciones en las mediciones.

\begin{table}[h]
    \centering
    \begin{tabular}{rr}
        \toprule
        \textbf{Cantidad de objetos ($n$)} & \textbf{Tiempo promedio (s)} \\
        \midrule
        10    & 0.000001689 \\
        20    & 0.000002922 \\
        50    & 0.000006530 \\
        90    & 0.000012225 \\
        100   & 0.000014572 \\
        500   & 0.000071820 \\
        1000  & 0.000148770 \\
        5000  & 0.000957233 \\
        10000 & 0.002141301 \\
        \bottomrule
    \end{tabular}
    \caption{Tiempos de ejecución experimentales del algoritmo Greedy.}
\end{table}

Para comparar los resultados experimentales con la complejidad teórica
$O(n \log n)$, se utilizó una curva de la forma

\[
T(n)=c\,n\log_2(n)
\]

donde la constante $c$ se ajustó a los tiempos experimentales mediante
mínimos cuadrados, obteniéndose aproximadamente

\[
c = 1.60 \times 10^{-8}.
\]

\begin{figure}[h]
    \centering
    \includegraphics[width=0.9\textwidth]{greedy_tiempos.png}
    \caption{Comparación entre los tiempos experimentales del algoritmo Greedy
    y la curva teórica ajustada $O(n \log n)$.}
    \label{fig:greedy_tiempos}
\end{figure}

Como puede observarse en la Figura \ref{fig:greedy_tiempos}, los tiempos
experimentales presentan un crecimiento muy cercano a la curva teórica
ajustada. Para los tamaños pequeños las diferencias son reducidas y pueden
estar influenciadas por el costo constante de ejecución y la precisión de la
medición. A medida que aumenta la cantidad de objetos, se observa con mayor
claridad el comportamiento esperado.

Los resultados experimentales son, por lo tanto, consistentes con el análisis
teórico realizado anteriormente, donde se determinó una complejidad temporal
de $O(n \log n)$, dominada por el ordenamiento de los objetos según su relación
beneficio/peso.

\newpage

% ==========================================
% PROBLEMA 3: PROGRAMACIÓN DINÁMICA
% ==========================================
\section{Problema 3: Programación dinámica}
\label{sec:problema3}

\subsection{Planteo tradicional (maximizar beneficio para capacidad fija)}

\subsubsection{Supuestos}
\textbf{Supuestos:} Se asume que la capacidad de la mochila $W$, así como los pesos $w_i$ y beneficios $v_i$ de cada objeto, son enteros no negativos.

\textbf{Ecuación de recurrencia:} Definimos $OPT(i, w)$ como el beneficio máximo que se puede obtener utilizando un subconjunto de los objetos del $1$ al $i$ con una capacidad máxima $w$.
\begin{align*}
    OPT(i, 0) &= 0 \quad \forall i \in \{0, \dots, n\} \\
    OPT(0, w) &= 0 \quad \forall w \in \{0, \dots, W\} \\
    OPT(i, w) &= OPT(i-1, w) \quad \text{si } w_i > w \\
    OPT(i, w) &= \max(OPT(i-1, w), OPT(i-1, w-w_i) + v_i) \quad \text{si } w_i \le w
\end{align*}

\subsubsection{Diseño, estructuras y pseudocódigo}
\textbf{Estructuras de datos utilizadas:}
\begin{itemize}
    \item \textbf{Lista de tuplas:} Para almacenar los objetos \texttt{(peso, beneficio)}.
    \item \textbf{Matriz:} Una tabla `dp` de dimensiones $(n+1) \times (W+1)$ para guardar en memoria los subproblemas ya resueltos.
\end{itemize}

\textbf{Pseudocódigo:}
\begin{lstlisting}[language=Python]
Funcion PD_Tradicional(capacidad, objetos):
    n = cantidad_objetos
    dp = matriz de (n+1) x (capacidad+1) inicializada en 0

    Para i desde 1 hasta n:
        peso_i, ben_i = objetos[i-1]
        Para w desde 0 hasta capacidad:
            Si peso_i <= w:
                dp[i][w] = max(dp[i-1][w], dp[i-1][w - peso_i] + ben_i)
            Sino:
                dp[i][w] = dp[i-1][w]

    # Reconstruccion de la solucion
    seleccionados = lista vacia
    w = capacidad
    Para i desde n hasta 1:
        Si dp[i][w] != dp[i-1][w]:      # el objeto i-1 fue incluido
            agregar objetos[i-1] a seleccionados
            w = w - peso_i

    Retornar seleccionados, dp[n][capacidad]
\end{lstlisting}

\subsubsection{Seguimiento}
Ejemplo reducido: Capacidad $W = 5$. Objetos: \texttt{O1(peso:3, ben:4)}, \texttt{O2(peso:2, ben:3)}.
\begin{table}[h!]
    \centering
    \begin{tabular}{@{}lcccccc@{}}
        \toprule
        \textbf{Objeto \textbackslash \ Capacidad} & \textbf{w=0} & \textbf{w=1} & \textbf{w=2} & \textbf{w=3} & \textbf{w=4} & \textbf{w=5} \\ \midrule
        $i=0$ (Sin objetos) & 0 & 0 & 0 & 0 & 0 & 0 \\
        $i=1$ ($O_1$: p=3, b=4) & 0 & 0 & 0 & \textbf{4} & 4 & 4 \\
        $i=2$ ($O_2$: p=2, b=3) & 0 & 0 & \textbf{3} & 4 & 4 & \textbf{7} \\ \bottomrule
    \end{tabular}
    \caption{Llenado de la matriz DP para el planteo tradicional.}
    \label{tab:pd_trad_seguimiento}
\end{table}
El beneficio óptimo es $dp[2][5] = 7$. Para reconstruir la solución se recorre la matriz de abajo hacia arriba partiendo de $w = 5$: $dp[2][5] = 7 \neq dp[1][5] = 4$, por lo que $O_2$ fue incluido y $w$ pasa a ser $3$; $dp[1][3] = 4 \neq dp[0][3] = 0$, por lo que $O_1$ también fue incluido. La solución es $\{O_1, O_2\}$ con peso $5$ y beneficio $7$.

\subsubsection{Complejidad temporal y espacial}
\begin{itemize}
    \item \textbf{Temporal:} $\mathcal{O}(n \cdot W)$. La inicialización de la matriz toma $\mathcal{O}(n \cdot W)$. El algoritmo consta de un doble bucle, donde el externo itera sobre los $n$ elementos y el interno sobre las capacidades desde $0$ hasta $W$. Como las operaciones dentro del bucle interno toman tiempo $\mathcal{O}(1)$, la complejidad temporal queda determinada por el llenado de la matriz, siendo $\mathcal{O}(n \cdot W)$. La reconstrucción de la solución recorre una vez las $n$ filas, lo que cuesta $\mathcal{O}(n)$ y no altera el orden.
    \item \textbf{Espacial:} $\mathcal{O}(n \cdot W)$. Se requiere alojar en memoria una matriz completa de tamaño $(n+1) \times (W+1)$ para almacenar los estados de la programación dinámica.
\end{itemize}

\bigskip

\subsection{Planteo alternativo (minimizar peso para beneficio fijo)}

\subsubsection{Supuestos}
\textbf{Supuestos:} Los pesos $w_i$ y los beneficios $v_i$ son enteros positivos. Definimos $V = \sum v_i$ como el beneficio que se obtendría incluyendo todos los objetos; es una cota superior del beneficio alcanzable y determina la cantidad de columnas de la matriz. El valor $\infty$ representa un beneficio inalcanzable con los objetos considerados.

\textbf{Ecuación de recurrencia:} Definimos $OPT(i, v)$ como el \textit{peso mínimo} de un subconjunto de los objetos del $1$ al $i$ cuyo beneficio es \textbf{al menos} $v$.
\begin{align*}
    OPT(i, 0) &= 0 \quad \forall i \in \{0, \dots, n\} \\
    OPT(0, v) &= \infty \quad \forall v \in \{1, \dots, V\} \\
    OPT(i, v) &= \min\big(OPT(i-1, v),\; OPT(i-1, \max(0, v - v_i)) + w_i\big) \quad \forall v \in \{1, \dots, V\}
\end{align*}
El primer término corresponde a no incluir el objeto $i$. El segundo lo incluye: si $v_i < v$, todavía faltan $v - v_i$ unidades de beneficio que deben cubrir los objetos anteriores; si $v_i \ge v$, el objeto por sí solo cubre el beneficio pedido y el resto del subconjunto puede ser vacío (columna $0$, de peso $0$). De ahí el $\max(0, v - v_i)$. Con esta definición la fila final es no decreciente en $v$, y la respuesta es el mayor $v$ tal que $OPT(n, v) \le W$: ese es el máximo beneficio que entra en la mochila.

Esta formulación difiere de la que exige obtener \textit{exactamente} un beneficio $v$ (en la que el objeto $i$ solo se considera si $v_i \le v$) únicamente en el contenido de las celdas; el beneficio óptimo que se obtiene es el mismo, pero elegimos describir la variante ``al menos $v$'' porque es la que implementa el código entregado.

\subsubsection{Diseño y pseudocódigo}
\textbf{Estructuras de datos utilizadas:}
\begin{itemize}
    \item \textbf{Lista de tuplas:} Para los objetos \texttt{(peso, beneficio)}.
    \item \textbf{Matriz:} Una tabla `dp` de $(n+1) \times (V+1)$ que inicializamos en infinito, salvo la columna $v=0$.
\end{itemize}

\textbf{Pseudocódigo:}
\begin{lstlisting}[language=Python]
Funcion PD_Alternativo(capacidad, objetos):
    n = cantidad_objetos
    Si n == 0:
        Retornar lista vacia, 0
    V = suma de todos los beneficios de los objetos
    dp = matriz de (n+1) x (V+1) inicializada en infinito

    Para i desde 0 hasta n:
        dp[i][0] = 0

    Para i desde 1 hasta n:
        peso_i, ben_i = objetos[i-1]
        Para v desde 0 hasta V:
            dp[i][v] = min(dp[i-1][v], dp[i-1][max(0, v - ben_i)] + peso_i)

    # Mayor beneficio cuyo peso minimo entra en la mochila
    mejor_v = 0
    Para v desde V hasta 0 (descendente):
        Si dp[n][v] <= capacidad:
            mejor_v = v
            Cortar

    # Reconstruccion de la solucion
    seleccionados = lista vacia
    v = mejor_v
    Para i desde n hasta 1:
        Si v <= 0:
            Cortar
        peso_i, ben_i = objetos[i-1]
        Si dp[i][v] != dp[i-1][v]:      # el objeto i-1 fue incluido
            agregar objetos[i-1] a seleccionados
            v = max(0, v - ben_i)

    Retornar seleccionados, mejor_v
\end{lstlisting}

\subsubsection{Seguimiento}
Ejemplo reducido: Capacidad $W = 5$. Objetos: \texttt{O1(peso:3, ben:4)}, \texttt{O2(peso:2, ben:3)}. $V = 7$.
\begin{table}[h!]
    \centering
    \begin{tabular}{@{}lcccccccc@{}}
        \toprule
        \textbf{Objeto \textbackslash \ Beneficio ($v$)} & \textbf{0} & \textbf{1} & \textbf{2} & \textbf{3} & \textbf{4} & \textbf{5} & \textbf{6} & \textbf{7} \\ \midrule
        $i=0$ (Sin obj.) & 0 & $\infty$ & $\infty$ & $\infty$ & $\infty$ & $\infty$ & $\infty$ & $\infty$ \\
        $i=1$ ($O_1$: p=3, b=4) & 0 & 3 & 3 & 3 & \textbf{3} & $\infty$ & $\infty$ & $\infty$ \\
        $i=2$ ($O_2$: p=2, b=3) & 0 & 2 & 2 & \textbf{2} & 3 & 5 & 5 & \textbf{5} \\ \bottomrule
    \end{tabular}
    \caption{Llenado de la matriz DP para el planteo alternativo: peso mínimo para obtener un beneficio de al menos $v$.}
    \label{tab:pd_alt_seguimiento}
\end{table}
En la fila $i=1$, con $O_1$ solo se puede alcanzar beneficio $4$ pagando peso $3$, y ese mismo peso cubre cualquier beneficio menor ($v \le 4$); los beneficios $5$ a $7$ son inalcanzables. En la fila $i=2$, para $v \le 3$ alcanza con $O_2$ (peso $2$); para $v = 4$ conviene $O_1$ solo (peso $3$) antes que $O_1 + O_2$ (peso $5$); y para $v \ge 5$ hacen falta ambos objetos (peso $5$).

La última fila es no decreciente, por lo que se la recorre desde $v = 7$ hacia abajo y se corta en el primer valor cuyo peso no supera la capacidad: $dp[2][7] = 5 \le 5$, de modo que el resultado es $7$. Para reconstruir la solución: $dp[2][7] = 5 \neq dp[1][7] = \infty$, así que $O_2$ fue incluido y $v$ pasa a ser $\max(0, 7 - 3) = 4$; $dp[1][4] = 3 \neq dp[0][4] = \infty$, así que $O_1$ también fue incluido. La solución es $\{O_1, O_2\}$ con peso $5$ y beneficio $7$, la misma que obtiene el planteo tradicional (Tabla~\ref{tab:pd_trad_seguimiento}).

\subsubsection{Complejidad temporal y espacial}
\begin{itemize}
    \item \textbf{Temporal:} $\mathcal{O}(n \cdot V)$. El cálculo del beneficio total $V$ toma tiempo $\mathcal{O}(n)$. Luego se construye la matriz con dimensiones $(n+1) \times (V+1)$. El llenado de la matriz consta de dos bucles anidados que realizan operaciones en tiempo $\mathcal{O}(1)$, por lo que demora $\mathcal{O}(n \cdot V)$. La búsqueda descendente sobre la última fila cuesta $\mathcal{O}(V)$ en el peor caso y la reconstrucción $\mathcal{O}(n)$. Por lo tanto, el tiempo total está dominado por el cálculo de la matriz y equivale a $\mathcal{O}(n \cdot V)$.
    \item \textbf{Espacial:} $\mathcal{O}(n \cdot V)$. La matriz `dp` debe almacenar estados para los $n$ objetos combinados con todos los beneficios posibles desde $0$ hasta $V$.
\end{itemize}

\subsection{Sets de datos utilizados}
Todos los sets se adjuntan en el directorio \texttt{pruebas/} y se midieron con el script \texttt{pruebas/grafico\_pd.py}, que ejecuta ambos algoritmos sobre cada archivo, verifica que devuelvan el mismo beneficio y genera las figuras de esta sección.
\begin{itemize}
    \item \textbf{Sets aleatorios} \texttt{mochila10.txt}, \texttt{mochila20.txt}, \texttt{mochila50.txt}, \texttt{mochila90.txt}, \texttt{mochila100.txt}, \texttt{mochila200.txt}, \texttt{mochila300.txt} y \texttt{mochila500.txt}, generados con el \texttt{crear\_mochila.py} de la cátedra: capacidad $W = 50n$, pesos uniformes en $[1, 200]$ y beneficios uniformes en $[1, 1000]$. El criterio fue cubrir dos órdenes de magnitud en $n$ con puntos intermedios suficientes para ajustar una curva. No se incluyó $n = 1000$ porque el planteo alternativo requiere una matriz de $1001 \times 500\,000$ celdas, que no entra en memoria con listas de Python.
    \item \textbf{Sets \textit{ad-hoc}} \texttt{mochila\_gana\_tradicional.txt} ($n = 50$, $W = 250$, beneficios del orden de $10^5$, $V = 14\,375\,970$) y \texttt{mochila\_gana\_alternativo.txt} ($n = 50$, $W = 12\,500\,000$, beneficios entre $1$ y $7$, $V = 244$). El criterio fue invertir la relación entre $W$ y $V$ para mostrar que el costo de cada algoritmo depende del parámetro que recorre su matriz y no solo de $n$.
\end{itemize}

Una observación que usamos en el ajuste: con la distribución de \texttt{crear\_mochila.py}, $W = 50n$ y el beneficio total esperado es $E[V] = 500{,}5 \cdot n$. Por lo tanto, sobre estos sets, $n \cdot W = 50n^2$ y $n \cdot V \approx 500n^2$: ambos algoritmos deberían verse cuadráticos en $n$, y el alternativo debería ser alrededor de $10$ veces más lento que el tradicional.

\subsection{Evaluación comparativa de ambos enfoques}

\subsubsection{Comportamiento con variación de \texorpdfstring{$W$}{W} vs. \texorpdfstring{$V$}{V}}
A la hora de comparar los tiempos empíricos entre ambos enfoques, se evidencia de forma clara el \textit{trade-off} que existe respecto a las variables $W$ (capacidad) y $V$ (beneficio total).

Para los sets generados con la distribución del enunciado, $V$ es alrededor de diez veces mayor que $W$. En las mediciones de la Tabla~\ref{tab:tabla_pd_tiempos} y la Figura~\ref{fig:pd_comparacion}, el \textbf{planteo tradicional es consistentemente más rápido}, con una relación entre tiempos que se mantiene entre $14$ y $17$ veces para todos los tamaños. Esa relación es algo mayor que el cociente $V/W \approx 10$ que predice el tamaño de las matrices, lo que atribuimos al costo adicional por celda del alternativo (aritmética con \texttt{float('inf')} y la llamada a \texttt{max}).

\begin{table}[H]
    \centering
    \begin{tabular}{rrrrrr}
        \toprule
        \textbf{$n$} & \textbf{$W$} & \textbf{$V$} & \textbf{PD Tradicional (s)} & \textbf{PD Alternativo (s)} & \textbf{Alt./Trad.} \\
        \midrule
        10  & 500   & 4665   & 0.000904 & 0.015621  & 17.3 \\
        20  & 1000  & 10842  & 0.004331 & 0.074841  & 17.3 \\
        50  & 2500  & 27537  & 0.028633 & 0.480360  & 16.8 \\
        90  & 4500  & 44637  & 0.094902 & 1.395703  & 14.7 \\
        100 & 5000  & 47445  & 0.117335 & 1.629541  & 13.9 \\
        200 & 10000 & 101222 & 0.478313 & 7.072966  & 14.8 \\
        300 & 15000 & 155699 & 1.116855 & 16.540967 & 14.8 \\
        500 & 25000 & 250073 & 3.228582 & 46.124868 & 14.3 \\
        \bottomrule
    \end{tabular}
    \caption{Tiempos de ejecución experimentales (mediana de varias corridas) para ambos algoritmos de programación dinámica en los sets aleatorios.}
    \label{tab:tabla_pd_tiempos}
\end{table}

\begin{figure}[H]
    \centering
    \includegraphics[width=0.85\textwidth]{pd_comparacion.png}
    \caption{Tiempo de ejecución de ambos planteos en función de $n$ sobre los sets aleatorios, en escala logarítmica en ambos ejes.}
    \label{fig:pd_comparacion}
\end{figure}

Para mostrar que la ventaja no es intrínseca del planteo tradicional sino de la relación entre $W$ y $V$, se midieron los dos sets \textit{ad-hoc} descritos en la sección anterior (Tabla~\ref{tab:pd_adhoc}):
\begin{itemize}
    \item \textbf{Caso favorable al tradicional:} capacidad ínfima ($W = 250$) y beneficios enormes ($V \approx 1{,}4 \cdot 10^7$). El tradicional llena una matriz de $51 \times 251$ celdas y responde en milisegundos; el alternativo necesita $51 \times 14\,375\,971$ celdas y demora minutos.
    \item \textbf{Caso favorable al alternativo:} capacidad colosal ($W = 12{,}5 \cdot 10^6$) y beneficios triviales ($V = 244$). Ahora es el tradicional el que recorre $51 \times 12\,500\,001$ celdas, mientras que el alternativo resuelve una matriz de $51 \times 245$ celdas.
\end{itemize}

\begin{table}[H]
    \centering
    \begin{tabular}{lrrrrr}
        \toprule
        \textbf{Set} & \textbf{$n$} & \textbf{$W$} & \textbf{$V$} & \textbf{PD Tradicional (s)} & \textbf{PD Alternativo (s)} \\
        \midrule
        \texttt{mochila\_gana\_tradicional.txt} & 50 & 250        & 14\,375\,970 & 0.0027   & 246.7900 \\
        \texttt{mochila\_gana\_alternativo.txt} & 50 & 12\,500\,000 & 244        & 137.7966 & 0.0043   \\
        \bottomrule
    \end{tabular}
    \caption{Tiempos de ejecución sobre los sets \textit{ad-hoc} (una corrida cada uno).}
    \label{tab:pd_adhoc}
\end{table}

En ambos casos el algoritmo cuya matriz es pequeña responde en milisegundos y el otro tarda minutos, con $n$ idéntico. Esto confirma que ninguno de los dos domina al otro: la elección depende de cuál de los dos parámetros, $W$ o $V$, es menor en la instancia a resolver.

\subsubsection{Curvas de ajuste teórico y medición real}
Para contrastar las mediciones con la complejidad calculada se ajustó, por mínimos cuadrados, una recta sin ordenada al origen $T = c \cdot (n \cdot W)$ para el tradicional y $T = c' \cdot (n \cdot V)$ para el alternativo, usando el $W$ y el $V$ reales de cada archivo. La constante resultante se interpreta como el tiempo que tarda cada algoritmo en calcular una celda de su matriz. Como sobre estos sets $W = 50n$ y $V \approx 500n$, las curvas ajustadas se ven como parábolas al graficarlas contra $n$ (Figuras~\ref{fig:pd_tradicional_ajuste} y~\ref{fig:pd_alternativo_ajuste}).

\begin{figure}[H]
    \centering
    \includegraphics[width=\textwidth]{pd_tradicional_ajuste.png}
    \caption{Planteo tradicional: tiempo medido frente al ajuste $c \cdot n \cdot W$, en escala lineal (izquierda) y logarítmica (derecha).}
    \label{fig:pd_tradicional_ajuste}
\end{figure}

\begin{figure}[H]
    \centering
    \includegraphics[width=\textwidth]{pd_alternativo_ajuste.png}
    \caption{Planteo alternativo: tiempo medido frente al ajuste $c' \cdot n \cdot V$, en escala lineal (izquierda) y logarítmica (derecha).}
    \label{fig:pd_alternativo_ajuste}
\end{figure}

Las constantes obtenidas fueron $c = 2{,}25 \cdot 10^{-7}$ s por celda para el tradicional y $c' = 3{,}49 \cdot 10^{-7}$ s por celda para el alternativo. Que el alternativo pague más por celda es coherente con lo observado en la Tabla~\ref{tab:tabla_pd_tiempos}: la relación de tiempos ($14$ a $17$) supera a la relación de tamaños de matriz ($V/W \approx 10$) justamente en el cociente $c'/c \approx 1{,}55$.

En el alternativo el ajuste reproduce todas las mediciones con un error relativo menor al $6\,\%$. En el tradicional el error es menor al $6\,\%$ entre $n = 20$ y $n = 200$, pero el ajuste sobreestima en un $25\,\%$ el punto $n = 10$ y subestima en un $13\,\%$ el punto $n = 500$. Lo primero se explica porque en $n = 10$ la matriz tiene apenas $5\,500$ celdas y el tiempo total (menos de un milisegundo) está dominado por costos fijos. Lo segundo lo atribuimos a que, al crecer $W$, los valores de la fila dejan de caber en la caché del procesador y el costo por celda aumenta levemente; en el alternativo este efecto es menor porque su costo por celda ya es mayor y lo enmascara. En ambos casos, la escala logarítmica muestra que la pendiente de las mediciones coincide con la de la curva teórica en todo el rango, lo que confirma el orden $\mathcal{O}(n \cdot W)$ y $\mathcal{O}(n \cdot V)$ calculado a partir del pseudocódigo.

Ambos enfoques resultan ser algoritmos pseudo-polinomiales exactos con fortalezas opuestas, lo que permite seleccionar el más adecuado basándose estrictamente en la magnitud de los parámetros de entrada: el tradicional cuando $W < V$, que es lo habitual en instancias con beneficios del orden de los pesos o mayores, y el alternativo cuando los beneficios son pequeños frente a la capacidad.

\newpage

% ==========================================
% PROBLEMA 4: PROGRAMACIÓN LINEAL
% ==========================================
\section{Problema 4: Programación lineal entera (PuLP)}
\label{sec:problema4}

\subsection{Supuestos}
\begin{itemize}
    \item \textbf{Datos enteros y positivos:} pesos, beneficios y capacidad son enteros positivos, tal como los genera \texttt{crear\_mochila.py}.
    \item \textbf{Solver externo:} el modelo se resuelve con CBC, el solver que viene incluido con PuLP (\texttt{PULP\_CBC\_CMD}). El tiempo medido incluye armar el modelo, pasárselo a CBC y leer la solución.
    \item \textbf{Solución exacta:} no se configura límite de tiempo ni tolerancia de \textit{gap}, por lo que CBC busca hasta probar que la solución es óptima. El resultado debe coincidir con el de los demás métodos exactos.
    \item \textbf{Lectura de la solución:} se asume que CBC devuelve las variables binarias con valores exactos $0$ o $1$, y se consideran seleccionados los objetos cuya variable vale $1$.
\end{itemize}

\subsection{Diseño}
\textbf{Modelo matemático:} sea $w_i$ el peso y $v_i$ el beneficio del objeto $i$, y $W$ la capacidad de la mochila. Se define una variable binaria $x_i$ por objeto, que vale $1$ si el objeto entra en la mochila y $0$ si no:
\begin{align*}
    \max \quad & \sum_{i=1}^{n} v_i x_i \\
    \text{s.a.} \quad & \sum_{i=1}^{n} w_i x_i \le W \\
    & x_i \in \{0, 1\} \quad \forall i \in \{1, \dots, n\}
\end{align*}

El modelo tiene $n$ variables y una sola restricción, sin importar el valor de $W$.

\textbf{Estructuras de datos utilizadas:}
\begin{itemize}
    \item \textbf{Lista de tuplas:} los objetos de entrada \texttt{(peso, beneficio)}.
    \item \textbf{Lista de variables binarias:} $n$ objetos \texttt{LpVariable}, uno por objeto.
    \item \textbf{Problema \texttt{LpProblem}:} guarda la función objetivo y la restricción de capacidad, cada una como una expresión lineal de $n$ términos.
    \item \textbf{Lista de índices:} los objetos seleccionados en la solución.
\end{itemize}

\textbf{Pseudocódigo:}
\begin{lstlisting}[language=Python]
Funcion MochilaLineal(capacidad, objetos):
    n = cantidad_objetos
    problema = nuevo problema de maximizacion

    Para i desde 0 hasta n - 1:
        x[i] = nueva variable binaria

    problema.objetivo = suma(beneficio_i * x[i])
    problema.agregar(suma(peso_i * x[i]) <= capacidad)

    Resolver problema con CBC

    seleccionados = [objetos[i] tal que x[i] == 1]
    valor_total = suma(beneficio de cada objeto en seleccionados)
    Retornar (seleccionados, valor_total)
\end{lstlisting}

\textbf{Resolución con CBC:} CBC usa \textit{branch and bound} (con planos de corte). En cada nodo resuelve la \textit{relajación lineal} del problema, que permite $0 \le x_i \le 1$. El valor de esa relajación es una cota superior del mejor beneficio alcanzable en ese nodo. Si la solución relajada es entera, es candidata a mejor solución. Si tiene una variable fraccionaria, se ramifica en dos subproblemas ($x_i = 0$ y $x_i = 1$). Se descarta todo nodo cuya cota no supere la mejor solución entera encontrada o que sea infactible.

\subsection{Seguimiento}
Ejemplo reducido: $W = 10$. Objetos: \texttt{O1(peso:5, ben:10)}, \texttt{O2(peso:4, ben:7)}, \texttt{O3(peso:6, ben:11)}. El modelo queda:
\[
    \max\ 10x_1 + 7x_2 + 11x_3 \quad \text{s.a.} \quad 5x_1 + 4x_2 + 6x_3 \le 10, \quad x_i \in \{0,1\}
\]

\begin{table}[h!]
    \centering
    \begin{tabular}{@{}lp{3cm}lrp{4.5cm}@{}}
        \toprule
        \textbf{Nodo} & \textbf{Restricciones} & \textbf{Relajación $(x_1,x_2,x_3)$} & \textbf{Cota} & \textbf{Acción} \\ \midrule
        Raíz & -- & $(1,\ 0,\ 5/6)$ & 19.17 & $x_3$ fraccionaria: ramifica. \\
        A & $x_3 = 0$ & $(1,\ 1,\ 0)$ & 17 & Entera. Mejor solución $= 17$. \\
        B & $x_3 = 1$ & $(0.8,\ 0,\ 1)$ & 19 & $19 > 17$ y $x_1$ fraccionaria: ramifica. \\
        B1 & $x_3 = 1,\ x_1 = 0$ & $(0,\ 1,\ 1)$ & 18 & Entera y $18 > 17$. Mejor solución $= 18$. \\
        B2 & $x_3 = 1,\ x_1 = 1$ & -- & -- & Infactible (peso $11 > 10$): se descarta. \\ \bottomrule
    \end{tabular}
\end{table}

La solución óptima es $x = (0, 1, 1)$, con beneficio $18$ y peso $10$, que coincide con la salida de \texttt{mochila\_lineal} (\texttt{([(4, 7), (6, 11)], 18)}). Tomar los objetos en orden de beneficio/peso daría O1 y O2, con beneficio $17$: la restricción de integralidad es la que obliga a ramificar. La tabla muestra la lógica del \textit{branch and bound}; internamente CBC aplica además preprocesamiento y cortes, y puede recorrer los nodos en otro orden.

\subsection{Complejidad}
\texttt{mochila\_lineal} tiene dos etapas con costos distintos.

\textbf{Armado del modelo y lectura de la solución: $O(n)$.} Se crea una variable por objeto, se escriben dos sumas de $n$ términos (objetivo y restricción) y al final se recorren las $n$ variables para ver cuáles valen $1$. La capacidad $W$ no influye en este costo. La memoria también es $O(n)$.

\textbf{Resolución con CBC: exponencial en el peor caso.} La mochila 0/1 es NP-difícil. En el peor caso el \textit{branch and bound} explora $O(2^n)$ nodos, y en cada uno resuelve una relajación lineal de costo polinomial (para una mochila con una sola restricción se resuelve ordenando por beneficio/peso, en $O(n \log n)$). Por lo tanto:
\[
    T(n) = O(n) + O(2^n \cdot n \log n) \quad \text{(peor caso)}
\]

\textbf{Comportamiento esperado en los sets de \texttt{crear\_mochila.py}.} La relajación lineal deja como mucho un objeto fraccionario. Como los pesos van de $1$ a $200$ y la capacidad es $50n$, cada objeto es chico frente a $W$: la cota de la relajación queda muy cerca del óptimo entero y la poda cierra la búsqueda en pocos nodos. Por eso en estos sets se espera que el tiempo crezca como $a \cdot n$ o $a \cdot n \log n$ (armado del modelo más el ordenamiento), sumado a un costo fijo: PuLP escribe el modelo en un archivo y ejecuta CBC como un proceso aparte, lo que domina el tiempo en instancias chicas. En otros sets (por ejemplo, con pesos grandes y parecidos a la capacidad), la búsqueda puede crecer como $2^n$.

\subsection{Resultados experimentales y ajuste curva teórica}
Los tiempos se obtuvieron con \texttt{pruebas/grafico\_lineal.py}. Para cada tamaño $n$ se generaron $5$ instancias con la misma distribución que \texttt{crear\_mochila.py} (pesos entre $1$ y $200$, beneficios entre $1$ y $1000$, capacidad $50n$), con semilla fija para que la medición sea reproducible. Cada instancia se resolvió una vez sin medir, para descartar el arranque en frío de CBC, y luego $3$ veces más. De esas $3$ corridas se toma la mediana. La tabla informa el promedio de las $5$ medianas y el rango entre la instancia más rápida y la más lenta. Los resultados por instancia quedan en \texttt{pruebas/lineal\_tiempos.csv}.

\begin{table}[h!]
    \centering
    \begin{tabular}{@{}rrrr@{}}
        \toprule
        \textbf{$n$} & \textbf{Promedio (s)} & \textbf{Mínimo (s)} & \textbf{Máximo (s)} \\ \midrule
        10    & 0.0255 & 0.0225 & 0.0333 \\
        20    & 0.0243 & 0.0226 & 0.0248 \\
        50    & 0.0269 & 0.0239 & 0.0336 \\
        90    & 0.0309 & 0.0252 & 0.0365 \\
        100   & 0.0354 & 0.0282 & 0.0453 \\
        500   & 0.0615 & 0.0458 & 0.0972 \\
        1000  & 0.0830 & 0.0477 & 0.1204 \\
        5000  & 0.2564 & 0.2128 & 0.3288 \\
        10000 & 0.5866 & 0.4146 & 0.7849 \\ \bottomrule
    \end{tabular}
    \caption{Tiempos de \texttt{mochila\_lineal} (5 instancias por tamaño).}
    \label{tab:pl_tiempos}
\end{table}

Además, sobre los sets del repositorio (\texttt{mochila10}, \texttt{20}, \texttt{50}, \texttt{90} y \texttt{1000}), el beneficio obtenido coincide con el de programación dinámica ($3569$, $9298$, $19402$, $36341$ y $411152$), y el peso de la solución nunca supera la capacidad.

\textbf{Ajuste a la curva teórica.} A partir del análisis de complejidad se ajustó $T(n) = a \cdot n \log_2 n + b$ por mínimos cuadrados. El término constante $b$ representa el costo fijo de escribir el modelo a un archivo y ejecutar CBC como proceso aparte. El ajuste minimiza el error relativo (cada punto pesa $1/T^2$): con mínimos cuadrados comunes los tamaños grandes dominan el resultado y el ajuste queda mal en los chicos. Se obtuvo:
\[
    a \approx 4.3 \times 10^{-6}\ \text{s}, \qquad b \approx 27\ \text{ms}
\]

\begin{figure}[h!]
    \centering
    \includegraphics[width=\textwidth]{lineal_tiempos.png}
    \caption{Tiempo de \texttt{mochila\_lineal} frente al ajuste $a \cdot n \log n + b$, en escala lineal (izquierda) y logarítmica (derecha). Las barras indican el rango entre instancias.}
    \label{fig:pl_tiempos}
\end{figure}

\subsection{Conclusiones}
\begin{itemize}
    \item \textbf{No se observa crecimiento exponencial.} Entre $n = 10$ y $n = 10000$ el tiempo pasa de unos $25$ ms a menos de $0.6$ s, y la curva sigue el ajuste $a \cdot n \log n + b$. Esto confirma lo anticipado en el análisis: en estos sets cada objeto es chico frente a la capacidad, la relajación lineal queda muy cerca del óptimo entero y CBC cierra el \textit{branch and bound} en pocos nodos. El peor caso $O(2^n)$ sigue siendo válido como cota, pero no aparece con esta distribución de datos.
    \item \textbf{En instancias chicas domina el costo fijo.} Hasta $n = 100$ el tiempo casi no cambia (entre $24$ y $35$ ms) y prácticamente todo es el valor de $b$: lanzar CBC y comunicarse con él por archivos. Por eso, para mochilas chicas, programación lineal es más lenta que métodos más simples, aunque su complejidad sea mejor en la práctica.
    \item \textbf{El tiempo depende de la instancia, no solo de $n$.} Para un mismo tamaño, el tiempo varía hasta un factor $2.5$ entre instancias (por ejemplo, entre $0.048$ y $0.120$ s para $n = 1000$). Esto se debe a la cantidad de nodos que explora CBC, que depende de los datos. Por eso se promediaron varias instancias por tamaño: con una sola, un caso desfavorable puede parecer un cambio de tendencia.
    \item \textbf{Distinguir $n$ de $n \log n$ no es concluyente.} Al duplicar $n$ de $5000$ a $10000$ el tiempo crece $\times 2.29$, más cerca de $n \log n$ ($\times 2.16$) que de $n$ ($\times 2$). Sin embargo, la dispersión entre instancias es mayor que esa diferencia, así que los datos son compatibles con ambas.
    \item \textbf{Solución exacta.} En todos los sets el resultado coincide con el de programación dinámica, como se espera de un método exacto.
\end{itemize}

\newpage

% ==========================================
% CONCLUSIÓN COMPARATIVA GENERAL
% ==========================================
\section{Conclusión comparativa global}
\label{sec:conclusion}

\subsection{Matriz resumen de algoritmos}

\begin{table}[h]
    \centering
    \begin{tabular}{lp{3.5cm}p{4cm}p{3.5cm}}
        \toprule
        \textbf{Algoritmo} & \textbf{Complejidad Teórica} & \textbf{Ventajas Principales} & \textbf{Limitaciones} \\
        \midrule
        Fuerza Bruta & $O(2^n)$ & Exacto, simple de formular & Intratable para $n > 30$ \\
        Backtracking & $O(2^n)$ peor caso & Poda ramas inviables & Depende del orden de poda \\
        Greedy ($1/2$) & $O(n \log n)$ & Extremadamente veloz & No garantiza óptimo global \\
        PD Tradicional & $O(n \cdot W)$ & Exacto & Sensible a capacidades $W$ gigantes \\
        PD Alternativa & $O(n \cdot V)$ & Exacto & Sensible a beneficios $V$ grandes \\
        Prog. Lineal & $O(2^n)$ peor caso (\textit{branch and bound}) & Exacto; modelo declarativo, fácil de extender con más restricciones & Costo fijo del solver; depende de bibliotecas externas \\
        \bottomrule
    \end{tabular}
    \caption{Comparación sintética de los 6 enfoques implementados.}
\end{table}

\subsection{Comparativa de tiempos en rangos comunes}
Gráfico general consolidado que enfrente los 6 algoritmos en el rango de tamaños donde todos fueron computacionalmente viables.

\subsection{Conclusiones finales}
Balance grupal sobre el compromiso entre tiempo de cálculo, consumo de memoria y exactitud de la solución obtenida.

\newpage

% ==========================================
% REFERENCIAS
% ==========================================
\section*{Referencias bibliográficas}
\addcontentsline{toc}{section}{Referencias Bibliográficas}

\begin{itemize}
    \item Cormen, T. H., Leiserson, C. E., Rivest, R. L., \& Stein, C. (2022). \textit{Introduction to Algorithms} (4ta ed.). MIT Press.
    \item Kleinberg, J., \& Tardos, É. (2006). \textit{Algorithm Design}. Pearson Education.
\end{itemize}

\end{document}